\subsection{Equal-data comparison with Chessformer}
\begin{table}[t]\centering\small
\caption{Equal-data comparison on the ALLIE 10k sample (EMA weights): top-1 (\%) / cross-entropy after $N$ training positions, same data stream and evaluation for all runs; throughput is end-to-end training positions per second on the 4-core CPU.}\label{tab:arch}
\begin{tabular}{lrccc}\toprule
Run & pos/s & 0.51M & 1.02M & 1.54M\\\midrule
Chessformer-5M (Maia-3 architecture) & 300 & 38.8 / 2.09 & 43.6 / 1.85 & 45.9 / 1.74 \\
Cloverb d256 L8 (no teacher) & 315 & 37.5 / 2.14 & 41.7 / 1.94 & 44.7 / 1.82 \\
Cloverb + online Maia-3-23M teacher & 153 & 42.9 / 1.90 & -- & -- \\
\bottomrule\end{tabular}\end{table}

We trained Cloverb (d=256, 8 layers, MLP$\times$2, FLOP-matched to Maia-3-5M: 0.30 vs.\ 0.32
GMAC per position) and the authors' own Chessformer-5M module (inputs exactly as in
\texttt{maia3.uci}) on the identical stream of 1.54M positions with the same optimiser, schedule
shape and evaluation (Table~\ref{tab:arch}). Two findings shaped the final system.

\emph{First, Chessformer's recipe is the bottleneck, not its architecture.} With a sensible
learning rate, pre-processing and kernels, the published 5M architecture learns fast:
45.9\% after 1.5M positions, i.e.\ 83\% of the way from chance to its published accuracy after
0.3\% of its published training data. The first Cloverb configuration learned \emph{more slowly
per position} (44.7\%) even though it trains 5\% faster per second; Section~\ref{sec:diag}
isolates why.

\emph{Second, a teacher is worth more than data.} Distilling the released Maia-3-23M online
(soft targets with weight 0.7 plus the human move) lifts Cloverb from 37.5\% to 42.9\% at 0.5M
positions, beyond what the student reaches with twice the data. Because the teacher's forward
pass costs about as much as a student training step, we amortise it: the final model trains for
several epochs on a fixed pool of positions labelled once by the clock-aware
Maia-3-23M-ponder checkpoint (top-16 moves, which hold 99.1\% of the probability mass).

\subsection{Diagnosing the per-sample gap}\label{sec:diag}
\begin{table}[t]\centering\small
\caption{Single-factor diagnostics after 512k training positions (EMA weights, ALLIE 10k sample; same data stream for every row). Indented rows each change one factor of the geometry-initialised model.}\label{tab:diag}
\begin{tabular}{lrrr}\toprule
Configuration & Top-1 & Top-3 & CE\\\midrule
Chessformer-5M (authors' module) & 38.8 & 63.4 & 2.088 \\
Cloverb, zero-initialised biases & 37.5 & 61.5 & 2.140 \\
\midrule
Cloverb + geometry initialisation & 44.4 & 71.1 & 1.823 \\
\quad + post-norm blocks & 45.0 & 72.4 & 1.791 \\
\quad + captured-piece type in history & 44.7 & 71.5 & 1.818 \\
\quad without dynamic templates ($T{=}0$) & 44.4 & 71.1 & 1.823 \\
\bottomrule\end{tabular}\end{table}


\subsection{Final model}
% main results table generated after the final run

